% Options for packages loaded elsewhere
% Options for packages loaded elsewhere
\PassOptionsToPackage{unicode}{hyperref}
\PassOptionsToPackage{hyphens}{url}
\PassOptionsToPackage{dvipsnames,svgnames,x11names}{xcolor}
%
\documentclass[
  letterpaper,
  DIV=11,
  numbers=noendperiod]{scrartcl}
\usepackage{xcolor}
\usepackage{amsmath,amssymb}
\setcounter{secnumdepth}{-\maxdimen} % remove section numbering
\usepackage{iftex}
\ifPDFTeX
  \usepackage[T1]{fontenc}
  \usepackage[utf8]{inputenc}
  \usepackage{textcomp} % provide euro and other symbols
\else % if luatex or xetex
  \usepackage{unicode-math} % this also loads fontspec
  \defaultfontfeatures{Scale=MatchLowercase}
  \defaultfontfeatures[\rmfamily]{Ligatures=TeX,Scale=1}
\fi
\usepackage{lmodern}
\ifPDFTeX\else
  % xetex/luatex font selection
\fi
% Use upquote if available, for straight quotes in verbatim environments
\IfFileExists{upquote.sty}{\usepackage{upquote}}{}
\IfFileExists{microtype.sty}{% use microtype if available
  \usepackage[]{microtype}
  \UseMicrotypeSet[protrusion]{basicmath} % disable protrusion for tt fonts
}{}
\makeatletter
\@ifundefined{KOMAClassName}{% if non-KOMA class
  \IfFileExists{parskip.sty}{%
    \usepackage{parskip}
  }{% else
    \setlength{\parindent}{0pt}
    \setlength{\parskip}{6pt plus 2pt minus 1pt}}
}{% if KOMA class
  \KOMAoptions{parskip=half}}
\makeatother
% Make \paragraph and \subparagraph free-standing
\makeatletter
\ifx\paragraph\undefined\else
  \let\oldparagraph\paragraph
  \renewcommand{\paragraph}{
    \@ifstar
      \xxxParagraphStar
      \xxxParagraphNoStar
  }
  \newcommand{\xxxParagraphStar}[1]{\oldparagraph*{#1}\mbox{}}
  \newcommand{\xxxParagraphNoStar}[1]{\oldparagraph{#1}\mbox{}}
\fi
\ifx\subparagraph\undefined\else
  \let\oldsubparagraph\subparagraph
  \renewcommand{\subparagraph}{
    \@ifstar
      \xxxSubParagraphStar
      \xxxSubParagraphNoStar
  }
  \newcommand{\xxxSubParagraphStar}[1]{\oldsubparagraph*{#1}\mbox{}}
  \newcommand{\xxxSubParagraphNoStar}[1]{\oldsubparagraph{#1}\mbox{}}
\fi
\makeatother


\usepackage{longtable,booktabs,array}
\usepackage{calc} % for calculating minipage widths
% Correct order of tables after \paragraph or \subparagraph
\usepackage{etoolbox}
\makeatletter
\patchcmd\longtable{\par}{\if@noskipsec\mbox{}\fi\par}{}{}
\makeatother
% Allow footnotes in longtable head/foot
\IfFileExists{footnotehyper.sty}{\usepackage{footnotehyper}}{\usepackage{footnote}}
\makesavenoteenv{longtable}
\usepackage{graphicx}
\makeatletter
\newsavebox\pandoc@box
\newcommand*\pandocbounded[1]{% scales image to fit in text height/width
  \sbox\pandoc@box{#1}%
  \Gscale@div\@tempa{\textheight}{\dimexpr\ht\pandoc@box+\dp\pandoc@box\relax}%
  \Gscale@div\@tempb{\linewidth}{\wd\pandoc@box}%
  \ifdim\@tempb\p@<\@tempa\p@\let\@tempa\@tempb\fi% select the smaller of both
  \ifdim\@tempa\p@<\p@\scalebox{\@tempa}{\usebox\pandoc@box}%
  \else\usebox{\pandoc@box}%
  \fi%
}
% Set default figure placement to htbp
\def\fps@figure{htbp}
\makeatother





\setlength{\emergencystretch}{3em} % prevent overfull lines

\providecommand{\tightlist}{%
  \setlength{\itemsep}{0pt}\setlength{\parskip}{0pt}}



 


\KOMAoption{captions}{tableheading}
\makeatletter
\@ifpackageloaded{caption}{}{\usepackage{caption}}
\AtBeginDocument{%
\ifdefined\contentsname
  \renewcommand*\contentsname{Table of contents}
\else
  \newcommand\contentsname{Table of contents}
\fi
\ifdefined\listfigurename
  \renewcommand*\listfigurename{List of Figures}
\else
  \newcommand\listfigurename{List of Figures}
\fi
\ifdefined\listtablename
  \renewcommand*\listtablename{List of Tables}
\else
  \newcommand\listtablename{List of Tables}
\fi
\ifdefined\figurename
  \renewcommand*\figurename{Figure}
\else
  \newcommand\figurename{Figure}
\fi
\ifdefined\tablename
  \renewcommand*\tablename{Table}
\else
  \newcommand\tablename{Table}
\fi
}
\@ifpackageloaded{float}{}{\usepackage{float}}
\floatstyle{ruled}
\@ifundefined{c@chapter}{\newfloat{codelisting}{h}{lop}}{\newfloat{codelisting}{h}{lop}[chapter]}
\floatname{codelisting}{Listing}
\newcommand*\listoflistings{\listof{codelisting}{List of Listings}}
\makeatother
\makeatletter
\makeatother
\makeatletter
\@ifpackageloaded{caption}{}{\usepackage{caption}}
\@ifpackageloaded{subcaption}{}{\usepackage{subcaption}}
\makeatother
\usepackage{bookmark}
\IfFileExists{xurl.sty}{\usepackage{xurl}}{} % add URL line breaks if available
\urlstyle{same}
\hypersetup{
  pdftitle={Group Assignment 1:Bayesian Reanalysis of the French Short Boredom Proneness Scale},
  pdfauthor={Nikola Sekulovski \& Maarten Marsman},
  colorlinks=true,
  linkcolor={blue},
  filecolor={Maroon},
  citecolor={Blue},
  urlcolor={Blue},
  pdfcreator={LaTeX via pandoc}}


\title{Group Assignment 1:Bayesian Reanalysis of the French Short
Boredom Proneness Scale}
\author{Nikola Sekulovski \& Maarten Marsman}
\date{}
\begin{document}
\maketitle


\section{Information about the
assignment}\label{information-about-the-assignment}

In this group assignment, you will \textbf{reanalyze an existing
dataset} using the Bayesian graphical modeling methodology introduced in
the workshop. You are provided with a short study description and a
\textbf{fully specified Methods section}. Your task is to:

\begin{itemize}
\tightlist
\item
  Implement the analysis exactly as implied by the Methods section
\item
  Make justified analytical decisions where choices are left open
\item
  Write a very bried \textbf{Results section} where the results are
  summarized.
\end{itemize}

You will work with the \textbf{French subsample (N = 490)} from the
original study, using \textbf{18 ordinal items} from the \textbf{Short
Boredom Proneness Scale (SBPS)} and the \textbf{Curiosity and
Exploration Inventory--II (CEI-II)}.

\newpage

\section{Introduction}\label{introduction}

Boredom proneness has been conceptualized as a trait-like tendency to
experience boredom frequently and intensely, and has been linked to a
broad range of psychological outcomes. Boredom proneness has been
conceptualized as a trait-like tendency to experience boredom frequently
and intensely, and has been linked to a broad range of psychological
outcomes. The Short Boredom Proneness Scale (SBPS) was developed as a
concise unidimensional measure of boredom proneness and has since been
translated into multiple languages.

Recent work has increasingly adopted a network perspective, in which
psychological constructs are modeled as systems of interacting
components rather than as manifestations of latent variables. Within
this framework, network models offer a flexible way to examine
conditional associations between individual questionnaire items while
controlling for all other variables in the system.

Therefore, in the original validation of the French version of the SBPS,
Martarelli et al.~(2022) examined the relationships between boredom
proneness, curiosity, and exploration using a regularized frequentist
network approach. Their analysis suggested positive associations between
boredom and curiosity items and negative associations between boredom
and exploration items.

The aim of the present study is to \textbf{reanalyze the French
subsample of the original dataset using Bayesian graphical modeling}. By
adopting a Bayesian approach, we explicitly quantify uncertainty in
network structure and parameters, and we distinguish between
\emph{absence of evidence} and \emph{evidence of absence} for network
edges. In particular, we estimate a Bayesian ordinal network model that
directly takes into account the ordinal measurement level of the
variables.

\section{Method}\label{method}

\subsection{Data}\label{data}

The dataset consists of responses from \textbf{N = 490 French-speaking
participants} (\(M_{age} = 33.8\); 73\% females and 27\% males). The
network includes 18 items, all measured on a \textbf{7-point Likert
scale}:

\begin{itemize}
\tightlist
\item
  8 items from the Short Boredom Proneness Scale (SBPS)
\item
  10 items from the Curiosity and Exploration Inventory--II (CEI-II)
\end{itemize}

Missing data was handled by listwise deletion.

\subsection{Model Specification}\label{model-specification}

To model conditional associations between items, we estimated a Bayesian
Ordinal Markov Random Field (OMRF, Marsman et al., 2025). This model
represents each item as a node in an undirected network, with edges
encoding conditional dependence relationships between item pairs while
controlling for all remaining variables in the network.

Bayesian model averaging was employed to integrate over uncertainty in
network structure. As a result, inference was not based on a single
selected network, but on the posterior distribution over possible
network structures.

\subsection{Prior Distributions}\label{prior-distributions}

\subsubsection{Structure Priors}\label{structure-priors}

The network structure was assigned a \textbf{Bernoulli prior} on edge
inclusion indicators. In the primary analysis, default prior settings
were used. To assess robustness, a \textbf{prior sensitivity analysis}
was conducted by repeating the analysis under alternative structure
priors reflecting different expectations about network density, ranging
from relatively sparse to relatively dense networks.

\subsubsection{Parameter Priors}\label{parameter-priors}

Conditional association parameters were assigned default priors that
regularize edge weights toward zero. To examine the influence of these
assumptions, a sensitivity analysis was performed by varying the scale
of the parameter priors, allowing for both more concentrated and more
diffuse prior distributions.

\subsection{Inference and Evidence
Quantification}\label{inference-and-evidence-quantification}

Evidence for the presence or absence of edges was quantified using
\textbf{Bayes factors}, comparing models in which a given edge was
included to models in which it was excluded. Bayes factors were used to
classify edges as showing evidence for inclusion, evidence for
exclusion, or inconclusive evidence.

Posterior summaries of edge parameters were obtained from the
model-averaged posterior distribution. Centrality measures were not
estimated or interpreted.

\section{Results}\label{results}

\subsection{Network Estimation}\label{network-estimation}

This section should report general characteristics of the estimated
network, including the number of nodes, the number of possible edges,
and an overall description of how evidence is distributed across
included, excluded, and inconclusive edges.

\subsection{Edge Evidence}\label{edge-evidence}

This section should describe the results of the Bayes factor analysis at
the edge level. Report how many edges show strong evidence for
inclusion, how many show strong evidence for exclusion, and how many
remain inconclusive. Particular attention should be paid to patterns of
associations within and between the SBPS and CEI-II items. Edge evidence
plots or summary tables should be referenced here.

\subsection{Edge Parameter Estimates}\label{edge-parameter-estimates}

This section should summarize the estimated edge parameters for edges
with non-negligible posterior support. Describe the direction and
relative magnitude of the associations, focusing on relationships
between boredom, curiosity, and exploration items. Interpretations
should remain descriptive and avoid causal language.

\subsection{Uncertainty in Parameter
Estimates}\label{uncertainty-in-parameter-estimates}

This section should discuss uncertainty in the estimated parameters.
Posterior intervals should be used to highlight which associations are
estimated with relatively high precision and which are more uncertain.

\subsection{Prior Sensitivity
Analysis}\label{prior-sensitivity-analysis}

This section should compare the primary analysis to the sensitivity
analyses. Discuss whether conclusions about edge inclusion, exclusion,
and parameter estimates are stable across different structure and
parameter priors, and highlight any results that are sensitive to prior
assumptions.

\section{References}\label{references}

\begin{itemize}
\item
  Martarelli, C. S., Baillifard, A., \& Audrin, C. (2022). A trait-based
  network perspective on the validation of the french short boredom
  proneness scale. \emph{European Journal of Psychological Assessment}.
  doi: https://doi.org/10.1027/1015-5759/a000718346
\item
  Marsman, M., van den Bergh, D., \& Haslbeck, J. M. B. (2025). Bayesian
  analysis of the ordinal Markov random field. \emph{Psychometrika, 90}
  , 146--182. doi: 10.1017/psy.2024.4
\end{itemize}




\end{document}
